\chapter*{مقدمة}
\addcontentsline{toc}{schapter}{مقدمة}
تُعدّ المعادلات التفاضلية من الركائز الأساسية في علم الرياضيات، لما لها من دورٍ محوري في تفسير الظواهر الطبيعية ونمذجتها بدقة، فضلاً عن تطبيقاتها الواسعة في مجالات متعددة مثل الفيزياء والهندسة والاقتصاد والعلوم الطبيعية. إذ تُمكّن هذه المعادلات من دراسة العلاقات بين المتغيرات المختلفة ومعدلات تغيرها، وهو ما يجعلها أداةً لا غنى عنها في فهم الأنظمة الديناميكية التي تتغير بمرور الزمن أو تبعًا لمؤثرات خارجية. ومن هذا المنطلق، فقد حظيت المعادلات التفاضلية باهتمامٍ كبير من قبل الباحثين، نظرًا لقدرتها على الربط بين الجانب النظري والتطبيقي في آنٍ واحد.\\
ومن بين الأنواع المهمة للمعادلات التفاضلية من الرتبة الأولى، تبرز معادلة برنولي ومعادلة ريكاتي لما لهما من أهمية نظرية وتطبيقية على حدٍ سواء. فهاتان المعادلتان تمثلان نماذج أساسية للمعادلات غير الخطية، والتي غالبًا ما تكون أكثر تعقيدًا من نظيراتها الخطية، الأمر الذي يدفع إلى البحث عن طرائق وأساليب فعّالة لتحليلها وحلها.\\
تتميز معادلة برنولي بإمكانية تحويلها إلى معادلة خطية من الرتبة الأولى بواسطة تغيير مناسب للمتغير، مما يجعل دراستها مثالاً مهماً على توظيف التحويلات الرياضية في تبسيط المسائل غير الخطية وتحويلها إلى صيغ يسهل التعامل معها. كما تُظهر هذه المعادلة كيف يمكن للأفكار البسيطة نسبيًا أن تُحدث فرقًا كبيرًا في حل مسائل تبدو معقدة للوهلة الأولى.\\
أما معادلة ريكاتي، فهي تُعد من المعادلات غير الخطية التي تحتل مكانة خاصة في التحليل الرياضي، نظرًا لارتباطها بعدد من الفروع المتقدمة في الرياضيات التطبيقية. إذ يمكن اختزالها إلى معادلة خطية إذا عُرف حل خاص لها، وهي خاصية تعكس عمق بنيتها الرياضية. كما تظهر هذه المعادلة في العديد من التطبيقات العملية، مثل ميكانيك الكم، ونظرية التحكم، وبعض مسائل الهندسة، مما يجعلها ذات أهمية كبيرة من الناحيتين النظرية والتطبيقية.\\
